\documentclass[10pt,a4paper]{article}
\usepackage[a4paper,margin=14mm,top=12mm,bottom=14mm]{geometry}
\usepackage{fontspec}
\setmainfont{Liberation Sans}
\usepackage{polyglossia}
\setdefaultlanguage{french}
\usepackage{amsmath}
\usepackage{xcolor}
\usepackage{graphicx}
\usepackage[most]{tcolorbox}
\usepackage{tabularx,booktabs,array,multirow}
\usepackage{enumitem}
\usepackage{fancyhdr}
\usepackage{lastpage}
\usepackage[hidelinks]{hyperref}
\definecolor{uimm}{HTML}{005677}
\definecolor{cyan}{HTML}{13A2DD}
\definecolor{fond}{HTML}{E6F2FA}
\definecolor{ok}{HTML}{2E7D32}
\definecolor{ko}{HTML}{B3261E}
\definecolor{violet}{HTML}{7B1FA2}
\setlength{\parindent}{0pt}
\setlength{\parskip}{3pt}
\setlist{nosep}
\setlist[itemize]{leftmargin=13pt,label=\textcolor{cyan}{\textbullet}}
\newcommand{\titre}[1]{\vspace{5pt}{\color{uimm}\large\bfseries #1}\par\vspace{-3pt}{\color{cyan}\rule{\linewidth}{1.2pt}}\par\vspace{0pt}}
\tcbset{boite/.style={colback=fond,colframe=uimm,boxrule=0.6pt,arc=2pt,left=6pt,right=6pt,top=4pt,bottom=4pt},
        defi/.style={colback=violet!6,colframe=violet,boxrule=0.8pt,arc=2pt,left=6pt,right=6pt,top=4pt,bottom=4pt,
                     title={\bfseries Défi},coltitle=white,colbacktitle=violet,fonttitle=\small}}
\newcommand{\bandeau}[3]{%
\begin{tcolorbox}[colback=uimm,colframe=uimm,arc=2pt,left=4pt,right=8pt,top=4pt,bottom=4pt]
\begin{minipage}[c]{18mm}\setlength{\fboxsep}{1.2pt}\fcolorbox{white}{white}{\includegraphics[height=16.5mm]{logo_uimm.jpg}}\end{minipage}\hspace{4mm}%
\begin{minipage}[c]{\dimexpr\linewidth-22mm\relax}\color{white}
{\small #1}\\[1pt]{\LARGE\bfseries #2}\\[2pt]{\large #3}\hfill{\small CQPM Technicien qualité}
\end{minipage}
\end{tcolorbox}}
\newcommand{\reponse}[1][1]{\par\vspace{2pt}{\color{gray!60}\foreach \i in {1,...,#1}{\rule{\linewidth}{0.3pt}\par\vspace{9pt}}}}
\usepackage{pgffor}
\pagestyle{fancy}\fancyhf{}
\renewcommand{\headrulewidth}{0pt}
\newcommand{\pied}[1]{\fancyfoot[L]{\footnotesize\color{uimm}#1}\fancyfoot[R]{\footnotesize\color{uimm}Page \thepage{} sur \pageref{LastPage}}}

\usepackage{tikz}
\usepackage{longtable}
\renewcommand{\bandeau}[3]{%
\begin{tcolorbox}[colback=uimm,colframe=uimm,arc=2pt,left=4pt,right=8pt,top=4pt,bottom=4pt]
\begin{minipage}[c]{18mm}\setlength{\fboxsep}{1.2pt}\fcolorbox{white}{white}{\includegraphics[height=16.5mm]{logo_uimm.jpg}}\end{minipage}\hspace{4mm}%
\begin{minipage}[c]{\dimexpr\linewidth-22mm\relax}\color{white}
{\small #1}\\[1pt]{\LARGE\bfseries #2}\\[2pt]{\large #3}
\end{minipage}
\end{tcolorbox}}
% en-tête d'une fiche : code, titre, question, bloc, durée, organisation, niveau
\newcommand{\entete}[7]{%
\begin{tcolorbox}[colback=uimm,colframe=uimm,arc=2pt,left=4pt,right=6pt,top=4pt,bottom=4pt]
\begin{minipage}[c]{18mm}\setlength{\fboxsep}{1.2pt}\fcolorbox{white}{white}{\includegraphics[height=16.5mm]{logo_uimm.jpg}}\end{minipage}\hspace{3mm}%
\begin{minipage}[c]{\dimexpr\linewidth-58mm\relax}\color{white}\raggedright
{\small Base SPC · #4}\\[1pt]{\Large\bfseries #1 · #2}\\[2pt]{\itshape #3}
\end{minipage}\hfill
\begin{minipage}[c]{33mm}\color{white}\small\raggedleft
Durée : \textbf{#5}\\ \textbf{#6}\\ Niveau #7 sur 3
\end{minipage}
\end{tcolorbox}}
\newcommand{\tagE}{{\setlength{\fboxsep}{1.5pt}\colorbox{cyan}{\color{white}\scriptsize\bfseries Ellistat}}}
\newcommand{\tagW}{{\setlength{\fboxsep}{1.5pt}\colorbox{uimm}{\color{white}\scriptsize\bfseries Page web}}}
\tcbset{pred/.style={colback=orange!6,colframe=orange!80!black,boxrule=0.6pt,arc=2pt,left=6pt,right=6pt,top=3pt,bottom=1pt,
                    title={\bfseries Je prédis, avant de commencer},coltitle=white,colbacktitle=orange!80!black,fonttitle=\small},
        retenir/.style={colback=fond,colframe=uimm,boxrule=0.6pt,arc=2pt,left=6pt,right=6pt,top=3pt,bottom=3pt,
                    title={\bfseries À retenir},coltitle=white,colbacktitle=uimm,fonttitle=\small},
        plus/.style={colback=violet!6,colframe=violet,boxrule=0.6pt,arc=2pt,left=6pt,right=6pt,top=3pt,bottom=3pt,
                    title={\bfseries Pour aller plus loin},coltitle=white,colbacktitle=violet,fonttitle=\small},
        act/.style={colback=uimm!10,colframe=uimm,boxrule=0.5pt,arc=2pt,left=5pt,right=5pt,top=2pt,bottom=2pt,before skip=7pt,after skip=3pt},
        verif/.style={colback=orange!6,colframe=orange!80!black,boxrule=0.6pt,arc=2pt,left=6pt,right=6pt,top=3pt,bottom=3pt,
                    title={\bfseries À vérifier sur poste avant la première séance},coltitle=white,colbacktitle=orange!80!black,fonttitle=\small}}
\newcommand{\fichier}[1]{\texttt{\small #1}}
\newcommand{\menu}[1]{\textsf{\textbf{#1}}}
\newcommand{\champ}[1]{\rule{#1}{0.3pt}}

\usepackage{colortbl}
\pied{Journée du 25/11/2026 · Trame de 6 h 30, matériel, copies, conseils · formateur}
\newcommand{\h}[2]{#1~h~#2}
\newcommand{\rr}{\cellcolor{orange!14}}
\begin{document}
\thispagestyle{fancy}
\bandeau{Journée du 25/11/2026 · 4 stagiaires · formateur}{Trame de 6 h 30 et réserves}{De 8 h 30 à 16 h 30 : le plan, le matériel, les copies, les conseils}
\vspace{2pt}
\begin{tcolorbox}[boite,left=5pt,right=5pt,top=2pt,bottom=2pt]\small
\textbf{Le temps réel.} Tes horaires donnent 7 h de présence, soit \textbf{6 h 30} sans les deux pauses de 15 min. Ton programme de 11 activités fait 7 h 20. J'y ajoute 8 min pour le quatrième opérateur de G2 (estimation de la fiche, non chronométrée), 10 min d'accueil et 15 min de bilan : \textbf{7 h 53}, à 7 min de tes 8 h minimum. Le plan retenu place \textbf{6 h 18} sur 6 h 30 : le tronc plus M9, avec 12 min de marge. M8 et G4 sont les réserves (lignes orangées). Tout faire demanderait d'aller 18 \% plus vite que les fiches.

\textbf{Un changement d'ordre à valider.} G1 passe \emph{après} le repas : le matin est celui du micromètre, l'après-midi celui du pied à coulisse et des pièces GUNT. G2 (R\&R à quatre opérateurs, jamais vu à l'écran) se trouve ainsi avant le repas avec 12 à 42 min de marge, et l'après-midi démarre par de la manipulation. Coût : G1 perd le lien direct avec M8, que le guide recommande de comparer. Avec G1 avant G2, le matin déborde de 3 min sur 12 h 30, sans M9.
\end{tcolorbox}
\par\vspace{4pt}
\renewcommand{\arraystretch}{1.1}\small
\begin{tabularx}{\linewidth}{|>{\raggedright\arraybackslash}p{15mm}|>{\bfseries}p{11mm}|>{\raggedright\arraybackslash}X|r|c|p{20mm}|}\hline
\rowcolor{uimm}\color{white}\textbf{Heure} & \color{white}\textbf{Fiche} & \color{white}\textbf{Activité} & \color{white}\textbf{Durée} & \color{white}\textbf{Rang} & \color{white}\textbf{Réel}\\\hline
\rowcolor{fond}\multicolumn{6}{|l|}{\textbf{\color{uimm}Matin · Le micromètre est-il capable ? (deux binômes A+B et C+D)}}\\ \hline
\h{8}{30} &  & Accueil : consignes, mémo Gage R\&R, ouvrir Ellistat & 10 min &  & \\ \hline
\h{8}{40} & M1 & La résolution d'abord & 15 min &  & \\ \hline
\h{8}{55} & M3 & Mon micromètre est-il capable ? & 40 min &  & \\ \hline
\h{9}{35} & M7 & Jusqu'où peut-il aller ? Aucune mesure nouvelle & 20 min &  & \\ \hline
\h{9}{55} & M16 & Deux grades, un micromètre (30 mesures, deux études) & 30 min &  & \\ \hline
\h{10}{25} &  & \textit{Pause} & 15 min &  & \\ \hline
\rowcolor{orange!14}\h{10}{40} & M9 & Deux opérateurs, une cale. Sans Ellistat & 30 min & R1 & \\ \hline
\rowcolor{orange!14}si avance & M8 & Juste partout ? Cinq cales, linéarité du micromètre & 45 min & R2 & \\ \hline
\h{11}{10} & G2 & Un micromètre, plusieurs tolérances. \textbf{Quatre opérateurs} : 60 min plus 8 & 68 min &  & \\ \hline
\h{12}{18} &  & Marge : enregistrer les projets Ellistat, ranger les cales & 12 min &  & \\ \hline
\rowcolor{gray!15}\h{12}{30} & \multicolumn{5}{l|}{\textit{Déjeuner jusqu'à 13 h 30}}\\ \hline
\rowcolor{fond}\multicolumn{6}{|l|}{\textbf{\color{uimm}Après-midi · Le pied à coulisse et les pièces GUNT}}\\ \hline
\h{13}{30} & G1 & Le pied à coulisse sur 150 mm & 45 min &  & \\ \hline
\h{14}{15} & G3 & Le nominal n'est pas le milieu (arbres PT105) & 45 min &  & \\ \hline
\h{15}{00} &  & \textit{Pause} & 15 min &  & \\ \hline
\rowcolor{orange!14}si avance & G4 & Bon ou rejet sur vingt plaques & 50 min & R3 & \\ \hline
\h{15}{15} & G5 & Quelle cote pose problème ? (plaques PT102) & 60 min &  & \\ \hline
\h{16}{15} &  & Bilan : un résultat et une décision par binôme & 15 min &  & \\ \hline
\h{16}{30} & \multicolumn{5}{l|}{\textit{Fin}}\\ \hline
\end{tabularx}
\par\smallskip
\footnotesize Rang : R1 est la dernière réserve à sauter, R3 la première. Les durées sont celles des fiches, non chronométrées avec ce groupe. Tu as noté que les stagiaires vont plus vite : écris l'heure réelle de fin de chaque activité dans la dernière colonne.

\titre{Points de contrôle et règles de bascule}
\renewcommand{\arraystretch}{1.1}\small
\begin{tabularx}{\linewidth}{|>{\raggedright\arraybackslash}p{30mm}|>{\raggedright\arraybackslash}p{40mm}|>{\raggedright\arraybackslash}X|}\hline
\rowcolor{uimm}\color{white}\textbf{Quand} & \color{white}\textbf{Test} & \color{white}\textbf{Décision}\\\hline
\h{10}{40}, fin de pause & Il est plus de \h{10}{52} & Saute M9. G2 doit démarrer à \h{11}{22} au plus tard pour finir à \h{12}{30}.\\ \hline
Fin de M9 (cible \h{11}{10}) & Il est \h{10}{37} ou plus tôt & Lance M8 (45 min) avant G2. Sinon, passe à G2.\\ \hline
Départ de G2 & Il est plus de \h{11}{22} & Fais les étapes Ellistat (5 et 6) sur un seul poste projeté, les autres regardent. Gain non mesuré.\\ \hline
\h{13}{30} & Aucun & G1 démarre à l'heure : le repas remet les compteurs à zéro.\\ \hline
Fin de G3 (cible \h{15}{00}) & Il est plus de \h{15}{05} & Réduis G5 à moins de cotes (étape 1 de la fiche). Garde 10 min de bilan au minimum.\\ \hline
Fin de G3 (cible \h{15}{00}) & Il est avant \h{14}{10} (très improbable) & G4 devient possible avant la pause, puis G5.\\ \hline
\end{tabularx}
\par\smallskip
\textbf{Place réelle de G4.} G4 demande 50 min et une référence que tu mesures au micromètre avant la séance. Elle n'a de place que de deux façons : à la place de G5, si tes pieds à coulisse ne mesurent proprement que trois cotes (voir « Résolution des pieds à coulisse » plus bas), ou si tu gagnes 50 min dans l'après-midi.

\textbf{Si un binôme finit avant l'autre.} M3, M9, M16 et G2 ont une case « Pour aller plus loin ». M7 est, selon le guide, une bonne activité d'attente.

\textbf{Si les stagiaires ne connaissent pas Ellistat} (question restée sans réponse) : les 10 min d'accueil deviennent E1 (« Mon premier projet »), E2 (« Saisir, coller, corriger », 15 min) passe avant M3, et tu paies ces 15 min en sautant M9.

\noindent\parbox{\linewidth}{\textbf{Au-delà du plan}, si tout va très vite : M10 (25 min, fabriquer son étalon, avant G1 pour le geste d'adhérence), M17 (70 min, R\&R sur dix cales à cote masquée), M15 (60 min, le micromètre dérive-t-il). Descriptions dans la trame de 8 h 20.}

\titre{Matériel à prévoir pour 4 stagiaires (deux binômes)}
\renewcommand{\arraystretch}{1.18}\small
\begin{tabularx}{\linewidth}{|>{\raggedright\arraybackslash}p{44mm}|>{\centering\arraybackslash}p{13mm}|>{\raggedright\arraybackslash}p{22mm}|>{\raggedright\arraybackslash}X|}\hline
\rowcolor{uimm}\color{white}\textbf{Matériel} & \color{white}\textbf{Quantité} & \color{white}\textbf{Activités} & \color{white}\textbf{Remarque}\\\hline
Poste avec Ellistat web, un compte par stagiaire & 4 + 1 & M3, M7, M8, M16, G1 à G5 & Le poste du formateur projette. Pas d'Ellistat en M1 et M9. Tester l'accès aux quatre pages web ci-dessous.\\ \hline
Micromètre 0-25 mm & 2 & M1, M3, M8, M9, M16 & Un par binôme. Zéro contrôlé, résolution relevée (0,01 ou 0,001), numéro noté (M3). Ton micromètre sert aussi à la référence de G4, mesurée avant la séance.\\ \hline
Micromètre 25-50 mm & 1 & G2 & Vérifier qu'il couvre le Ø 48 et relever sa résolution (le jeu de secours suppose 0,01).\\ \hline
Pied à coulisse 150 mm & 4 & M1, G1, G3, G4, G5 & M1 : un au 1/50 par binôme. G1 : celui de la mallette PT102. G4 : un par contrôleur, donc 3. G5 : avec jauge de profondeur. Relever la résolution de chacun.\\ \hline
Réglet & 2 & M1 & Ne figure pas dans la liste de ton outil.\\ \hline
Coffret de cales, classe 1 & 2 & M3, M8, M9, M16, G1 & Un par binôme. Cales utilisées : 10 mm (M3, M9, M16) ; 5, 10, 15, 20, 25 mm (M8) ; 10, 40, 50, 60, 90, 100 mm (G1). À vérifier dans le coffret.\\ \hline
Cale céramique 10 mm, grade 0 & 2 & M16 & Une par binôme. Avec une seule, les binômes passent l'un après l'autre. Repli du guide : deux cales acier de grades différents, ou deux piles de même valeur.\\ \hline
Plaques PT102 numérotées 1 à 20 & 20 & G1, G4, G5 & Deux mallettes. G1 n'en prend qu'une par binôme.\\ \hline
Arbres PT105 numérotés 1 à 20 & 20 & G2, G3 & Deux mallettes. G2 en prend 10 et demande un repère au feutre sur le Ø 48.\\ \hline
Petit matériel & 4 de chaque & toutes & Calculatrices ; feuilles quadrillées (biais en fonction de la longueur en M8, Pareto en G5) ; un cache pour les valeurs (M9) ; chiffons non pelucheux (M16, G1, G2 : les fiches demandent de nettoyer) ; un feutre fin (G2).\\ \hline
Tableau blanc ou paperboard & 1 & M8, G5 & G5 : les deux classements côte à côte, au tableau.\\ \hline
\end{tabularx}
\par\smallskip
\textbf{Pages web} (adresses complètes, comme sur les fiches).
\begin{itemize}[leftmargin=14pt,itemsep=1pt]\footnotesize
\item M9 et M16 : \texttt{https://sjaubert.github.io/000\_OI/SPC/Pages/Comparer\_deux\_series.html}
\item G2 : \texttt{https://sjaubert.github.io/000\_OI/SPC/Pages/L\_operateur\_compte-t-il.html}
\item M3, pour aller plus loin : \texttt{https://sjaubert.github.io/000\_OI/SPC/Pages/L\_instrument\_est-il\_capable.html}
\item Jeux de secours : \texttt{https://sjaubert.github.io/000\_OI/SPC/Pages/Donnees\_a\_copier.html} avec les ancres \texttt{\#instr} (M3, M7), \texttt{\#lin} (M8), \texttt{\#operateurs} (M9), \texttt{\#grades} (M16), \texttt{\#lin7} (G1), \texttt{\#rrarbres} et \texttt{\#rrarbres4} (G2), \texttt{\#decentre} (G3), \texttt{\#tri20} (G4), \texttt{\#cotes7} (G5).
\end{itemize}

\titre{Photocopies pour 4 stagiaires}
\renewcommand{\arraystretch}{1.18}\small
\begin{tabularx}{\linewidth}{|>{\bfseries}p{14mm}|>{\raggedright\arraybackslash}X|c|c|c|}\hline
\rowcolor{uimm}\color{white}\textbf{Fiche} & \color{white}\textbf{Contenu des pages} & \color{white}\textbf{Pages} & \color{white}\textbf{Copies} & \color{white}\textbf{Total}\\\hline
Mémo & Le Gage R\&R, c'est quoi ? (lu à l'accueil) & 1 & 4 & 4\\ \hline
M1 & Résolution, tableau & 1 & 4 & 4\\ \hline
M3 & Grille de 25 mesures, Ellistat & 1 & 4 & 4\\ \hline
M7 & Tableau Cg et Cgk & 1 & 4 & 4\\ \hline
M16 & Fiche et grille de 15 mesures, puis questions & 2 & 4 & 8\\ \hline
\rr M9 & Grilles de 15 mesures (R1) & \rr 1 & \rr 4 & \rr 4\\ \hline
G2 & Fiche, questions, plan PT105 & 3 & 4 & 12\\ \hline
G1 & Fiche, questions et encadré ISO 2768, plan PT102 & 3 & 4 & 12\\ \hline
G3 & Fiche, questions, plan PT105 & 3 & 4 & 12\\ \hline
G5 & Fiche, questions, plan PT102 & 3 & 4 & 12\\ \hline
\rowcolor{fond}\multicolumn{2}{|l|}{\textbf{Plan retenu}} & \textbf{19} & & \textbf{76}\\ \hline
\rr M8 & Cinq cales, tableau (R2) & \rr 1 & \rr 4 & \rr 4\\ \hline
\rr G4 & Fiche, plan PT102 avec encadré (R3) & \rr 2 & \rr 4 & \rr 8\\ \hline
\rowcolor{fond}\multicolumn{2}{|l|}{\textbf{Total pour les stagiaires}} & \textbf{22} & & \textbf{88}\\ \hline
\end{tabularx}
\par\smallskip
\textbf{À tirer en tout : 141 pages A4.} 88 pour les stagiaires, 22 pour ton jeu complet, 22 pour un jeu de rechange, 9 pour les pages du guide qui portent les corrigés : 8 à 10, 12 et 13, 17 à 20.

Imprime en recto simple et n'agrafe pas la dernière page des fiches G : le plan se pose à côté des pages de travail. Option économe : remplace la page 3 de G1, G2, G3 et G5 (16 pages) par deux exemplaires A3 de chaque plan (\texttt{plan\_plaque\_PT102.pdf}, \texttt{plan\_arbre\_PT105.pdf}, dans \texttt{SPC\textbackslash Sources\textbackslash fiches}). La page de plan de G4 porte l'encadré ISO 2768 : elle s'imprime.

\newpage
\titre{Articuler les activités : le fil de la journée}
\small
\textbf{La question qui court toute la journée :} peut-on se fier à ce qu'on mesure ? Le matin, on juge le micromètre : assez fin (M1), juste (M3, M7), jugé sur une bonne cale (M16), reproductible d'une personne à l'autre (M9, G2). L'après-midi, on juge un instrument plus grossier, le pied à coulisse (G1), puis les pièces qu'il mesure (G3, G5). Chaque activité répond à la question que la précédente laisse ouverte : dis-la à voix haute à chaque passage.
\par\vspace{3pt}
\renewcommand{\arraystretch}{1.2}\footnotesize
\begin{tabularx}{\linewidth}{|>{\bfseries\raggedright\arraybackslash}p{19mm}|>{\raggedright\arraybackslash}p{56mm}|>{\raggedright\arraybackslash}X|}\hline
\rowcolor{uimm}\color{white}\textbf{Passage} & \color{white}\textbf{Ce que tu dis} & \color{white}\textbf{Ce que tu surveilles}\\\hline
M1 vers M3 & « Le micromètre est assez fin pour 0,15 mm. Est-il juste ? » & Faire toucher les instruments avant de calculer (M1). La cale de 10 mm n'est pas au milieu de la tolérance de la butée (11,925) : le dire (guide M3). Le secrétaire saisit au fur et à mesure : 25 mesures prennent 10 à 15 min.\\ \hline
M3 vers M7 & « Même étude. On ne mesure plus rien : on ne change que la tolérance. » & Les stagiaires ne ferment ni ne suppriment l'étude de M3. Le changement de LSI et LSS dans une étude existante n'a pas été vu à l'écran : essaie-le avant.\\ \hline
M7 vers M16 & « Tout repose sur la cale. Une cale de grade 0 change-t-elle le résultat ? » & Prédiction écrite avant de commencer. L'activité est décevante par construction : ne conclus pas que les deux cales sont identiques (guide M16).\\ \hline
M16 vers M9 & « Même page web, mais on compare deux personnes au lieu de deux cales. » & Personne n'a « raison » : la cale dit qui est le plus juste. Cache sur les valeurs de l'autre (guide M9).\\ \hline
M9 vers G2 (M8 entre les deux si avance) & « Deux personnes ne trouvent pas la même moyenne. L'étude R\&R va mesurer cela sur dix arbres. » & G2 : ne pas annoncer la suite du tableau avant la prédiction ; seul T change entre les quatre lignes. Ressaisir les tolérances dans « Voir Plus » si Ellistat ne les garde pas ; à quatre opérateurs, contrôler qu'il affiche quatre tableaux. M8 : un mesure, un note, échange à mi-parcours ; tracer le biais avant d'ouvrir Ellistat ; saisir 5 produits (Ellistat propose 4).\\ \hline
G2 vers G1 (repas) & « Ce matin le micromètre. Cet après-midi, un instrument moins fin, sur les pièces des mallettes. » & Montre le geste d'adhérence des cales avant G1 (la fiche renvoie à M10) ; ne jamais séparer deux cales en tirant. Fais prédire le sens de l'erreur. Contrôle que la mesure 1 n'est pas égale à la valeur vraie dans les sept blocs (piège du 19/11).\\ \hline
G1 vers G3 & « On sait ce que vaut le pied à coulisse. On mesure maintenant des pièces : la cote est-elle centrée ? » & Le premier tableau (limites, milieu, T) se remplit avant de mesurer. La mauvaise lecture 30,3 ± 0,4 donne un Ppk rassurant à tort : insiste.\\ \hline
G3 vers G5 & « Sur deux cotes d'un arbre, le Ppk a parlé de centrage. Sur sept cotes d'une plaque, laquelle corriger d'abord ? » & 140 mesures : deux binômes, quelques cotes chacun. Essaie chaque cote sur une plaque avant. Mets les deux classements côte à côte au tableau.\\ \hline
G5 vers bilan & « Un instrument, une tolérance pour laquelle il est capable, une cote que tu corrigerais en premier. » & Pas de notion nouvelle dans les 15 dernières minutes.\\ \hline
\end{tabularx}
\par\smallskip
\normalsize
\textbf{Le rythme de chaque activité} (il tient dans toutes les fiches) : 1. la prédiction écrite, sans annoncer le résultat ; 2. la mesure, un mesure et un note, puis on échange ; 3. la saisie, chaque stagiaire dans son projet Ellistat ; 4. deux ou trois chiffres lus à l'écran ; 5. un stagiaire lit « À retenir » à voix haute, et tu notes l'heure réelle dans la trame.

\textbf{Organisation.} Binômes fixes toute la journée, sauf G2 (les quatre sont opérateurs, chacun note pour le suivant) et G4 (trois contrôleurs, un secrétaire). Les chiffres des corrigés viennent de jeux de secours simulés : ne les présente jamais comme le résultat attendu avec les vraies pièces.

\titre{Résolution des pieds à coulisse : ce qu'elle change}
Règle de M1 : résolution d'au plus 5 \% de T, l'intervalle de tolérance. Calcul pour les cotes de la journée :
\par\vspace{2pt}
\renewcommand{\arraystretch}{1.15}\footnotesize
\begin{tabularx}{\linewidth}{|>{\raggedright\arraybackslash}X|c|c|c|}\hline
\rowcolor{uimm}\color{white}\textbf{Cotes} & \color{white}\textbf{T (mm)} & \color{white}\textbf{Résolution maximale} & \color{white}\textbf{Pied au 1/50 (0,02)}\\\hline
G1 et G5 : a (±0,3) & 0,6 & 0,03 & convient\\ \hline
G3 : b, d ; G5 : m, n & 0,4 & 0,02 & à la limite (5 \%)\\ \hline
G4 : q ; G5 : d, q, r & 0,3 & 0,015 & ne convient pas (6,7 \%)\\ \hline
G1 et G5 : o (±0,1) & 0,2 & 0,01 & ne convient pas (10 \%)\\ \hline
\end{tabularx}
\par\smallskip
\normalsize
Avec un pied au 1/100, toutes les cotes passent. Avec un pied au 1/50, G5 ne garde que a, m et n : décide avant la séance entre G5 (courte, Pareto probablement plat) et G4. Le jeu de secours de G4 suppose des lectures à 0,05 mm.

\titre{À vérifier avant le 25/11}
\begin{itemize}[leftmargin=14pt,itemsep=1pt]\small
\item \textbf{Résolution réelle} de chaque micromètre et pied à coulisse : elle décide des verdicts de M1, G1, G2, G4, G5.
\item \textbf{Ellistat à quatre opérateurs (G2)} : jamais vu à l'écran. Écrans non vus avec ces données : G2 à G5 ; G3 (Détail avec tolérances vu pour E4 seulement) ; G5 (plusieurs colonnes : une à la fois) ; M7 (modifier LSI et LSS d'une étude existante).
\item \textbf{Pièces réelles} : les arbres et plaques sont probablement proches du milieu de leurs tolérances. Les résultats de G3 et G5 peuvent être plats (guide). Pour G4, fixe des limites d'exercice resserrées de façon que 6 à 8 plaques soient à moins de 0,03 mm d'une limite.
\item \textbf{Convention GUNT} pour b (30,3 +0,4) et d (55,2 −0,4) : « l'autre écart vaut zéro » n'est pas confirmé par GUNT.
\item \textbf{Coffrets} : présence des cales listées ; grade et certificat de la céramique de 10 mm ; cales et micromètres dans la salle une heure avant, zéro contrôlé.
\item \textbf{Durées} : toutes sont des estimations. M16 est serrée (30 mesures, deux études Ellistat, une page web).
\end{itemize}
\end{document}
